% Options for packages loaded elsewhere
% Options for packages loaded elsewhere
\PassOptionsToPackage{unicode}{hyperref}
\PassOptionsToPackage{hyphens}{url}
%
\documentclass[
  ignorenonframetext,
  aspectratio=169,
]{beamer}
\newif\ifbibliography
\usepackage{pgfpages}
\setbeamertemplate{caption}[numbered]
\setbeamertemplate{caption label separator}{: }
\setbeamercolor{caption name}{fg=normal text.fg}
\beamertemplatenavigationsymbolsempty
% remove section numbering
\setbeamertemplate{part page}{
  \centering
  \begin{beamercolorbox}[sep=16pt,center]{part title}
    \usebeamerfont{part title}\insertpart\par
  \end{beamercolorbox}
}
\setbeamertemplate{section page}{
  \centering
  \begin{beamercolorbox}[sep=12pt,center]{section title}
    \usebeamerfont{section title}\insertsection\par
  \end{beamercolorbox}
}
\setbeamertemplate{subsection page}{
  \centering
  \begin{beamercolorbox}[sep=8pt,center]{subsection title}
    \usebeamerfont{subsection title}\insertsubsection\par
  \end{beamercolorbox}
}
% Prevent slide breaks in the middle of a paragraph
\widowpenalties 1 10000
\raggedbottom
\AtBeginPart{
  \frame{\partpage}
}
\AtBeginSection{
  \ifbibliography
  \else
    \frame{\sectionpage}
  \fi
}
\AtBeginSubsection{
  \frame{\subsectionpage}
}
\usepackage{iftex}
\ifPDFTeX
  \usepackage[T1]{fontenc}
  \usepackage[utf8]{inputenc}
  \usepackage{textcomp} % provide euro and other symbols
\else % if luatex or xetex
  \usepackage{unicode-math} % this also loads fontspec
  \defaultfontfeatures{Scale=MatchLowercase}
  \defaultfontfeatures[\rmfamily]{Ligatures=TeX,Scale=1}
\fi
\usepackage{lmodern}

\usetheme[]{Boadilla}
\usecolortheme[]{default}
\usefonttheme[]{professionalfonts}
\ifPDFTeX\else
  % xetex/luatex font selection
\fi
% Use upquote if available, for straight quotes in verbatim environments
\IfFileExists{upquote.sty}{\usepackage{upquote}}{}
\IfFileExists{microtype.sty}{% use microtype if available
  \usepackage[]{microtype}
  \UseMicrotypeSet[protrusion]{basicmath} % disable protrusion for tt fonts
}{}
\makeatletter
\@ifundefined{KOMAClassName}{% if non-KOMA class
  \IfFileExists{parskip.sty}{%
    \usepackage{parskip}
  }{% else
    \setlength{\parindent}{0pt}
    \setlength{\parskip}{6pt plus 2pt minus 1pt}}
}{% if KOMA class
  \KOMAoptions{parskip=half}}
\makeatother

\usepackage{color}
\usepackage{fancyvrb}
\newcommand{\VerbBar}{|}
\newcommand{\VERB}{\Verb[commandchars=\\\{\}]}
\DefineVerbatimEnvironment{Highlighting}{Verbatim}{commandchars=\\\{\}}
% Add ',fontsize=\small' for more characters per line
\usepackage{framed}
\definecolor{shadecolor}{RGB}{241,243,245}
\newenvironment{Shaded}{\begin{snugshade}}{\end{snugshade}}
\newcommand{\AlertTok}[1]{\textcolor[rgb]{0.68,0.00,0.00}{#1}}
\newcommand{\AnnotationTok}[1]{\textcolor[rgb]{0.37,0.37,0.37}{#1}}
\newcommand{\AttributeTok}[1]{\textcolor[rgb]{0.40,0.45,0.13}{#1}}
\newcommand{\BaseNTok}[1]{\textcolor[rgb]{0.68,0.00,0.00}{#1}}
\newcommand{\BuiltInTok}[1]{\textcolor[rgb]{0.00,0.23,0.31}{#1}}
\newcommand{\CharTok}[1]{\textcolor[rgb]{0.13,0.47,0.30}{#1}}
\newcommand{\CommentTok}[1]{\textcolor[rgb]{0.37,0.37,0.37}{#1}}
\newcommand{\CommentVarTok}[1]{\textcolor[rgb]{0.37,0.37,0.37}{\textit{#1}}}
\newcommand{\ConstantTok}[1]{\textcolor[rgb]{0.56,0.35,0.01}{#1}}
\newcommand{\ControlFlowTok}[1]{\textcolor[rgb]{0.00,0.23,0.31}{\textbf{#1}}}
\newcommand{\DataTypeTok}[1]{\textcolor[rgb]{0.68,0.00,0.00}{#1}}
\newcommand{\DecValTok}[1]{\textcolor[rgb]{0.68,0.00,0.00}{#1}}
\newcommand{\DocumentationTok}[1]{\textcolor[rgb]{0.37,0.37,0.37}{\textit{#1}}}
\newcommand{\ErrorTok}[1]{\textcolor[rgb]{0.68,0.00,0.00}{#1}}
\newcommand{\ExtensionTok}[1]{\textcolor[rgb]{0.00,0.23,0.31}{#1}}
\newcommand{\FloatTok}[1]{\textcolor[rgb]{0.68,0.00,0.00}{#1}}
\newcommand{\FunctionTok}[1]{\textcolor[rgb]{0.28,0.35,0.67}{#1}}
\newcommand{\ImportTok}[1]{\textcolor[rgb]{0.00,0.46,0.62}{#1}}
\newcommand{\InformationTok}[1]{\textcolor[rgb]{0.37,0.37,0.37}{#1}}
\newcommand{\KeywordTok}[1]{\textcolor[rgb]{0.00,0.23,0.31}{\textbf{#1}}}
\newcommand{\NormalTok}[1]{\textcolor[rgb]{0.00,0.23,0.31}{#1}}
\newcommand{\OperatorTok}[1]{\textcolor[rgb]{0.37,0.37,0.37}{#1}}
\newcommand{\OtherTok}[1]{\textcolor[rgb]{0.00,0.23,0.31}{#1}}
\newcommand{\PreprocessorTok}[1]{\textcolor[rgb]{0.68,0.00,0.00}{#1}}
\newcommand{\RegionMarkerTok}[1]{\textcolor[rgb]{0.00,0.23,0.31}{#1}}
\newcommand{\SpecialCharTok}[1]{\textcolor[rgb]{0.37,0.37,0.37}{#1}}
\newcommand{\SpecialStringTok}[1]{\textcolor[rgb]{0.13,0.47,0.30}{#1}}
\newcommand{\StringTok}[1]{\textcolor[rgb]{0.13,0.47,0.30}{#1}}
\newcommand{\VariableTok}[1]{\textcolor[rgb]{0.07,0.07,0.07}{#1}}
\newcommand{\VerbatimStringTok}[1]{\textcolor[rgb]{0.13,0.47,0.30}{#1}}
\newcommand{\WarningTok}[1]{\textcolor[rgb]{0.37,0.37,0.37}{\textit{#1}}}

\usepackage{longtable,booktabs,array}
\newcounter{none} % for unnumbered tables
\usepackage{calc} % for calculating minipage widths
\usepackage{caption}
% Make caption package work with longtable
\makeatletter
\def\fnum@table{\tablename~\thetable}
\makeatother
\usepackage{graphicx}
\makeatletter
\newsavebox\pandoc@box
\newcommand*\pandocbounded[1]{% scales image to fit in text height/width
  \sbox\pandoc@box{#1}%
  \Gscale@div\@tempa{\textheight}{\dimexpr\ht\pandoc@box+\dp\pandoc@box\relax}%
  \Gscale@div\@tempb{\linewidth}{\wd\pandoc@box}%
  \ifdim\@tempb\p@<\@tempa\p@\let\@tempa\@tempb\fi% select the smaller of both
  \ifdim\@tempa\p@<\p@\scalebox{\@tempa}{\usebox\pandoc@box}%
  \else\usebox{\pandoc@box}%
  \fi%
}
% Set default figure placement to htbp
\def\fps@figure{htbp}
\makeatother





\setlength{\emergencystretch}{3em} % prevent overfull lines

\providecommand{\tightlist}{%
  \setlength{\itemsep}{0pt}\setlength{\parskip}{0pt}}



 


\makeatletter
\@ifpackageloaded{caption}{}{\usepackage{caption}}
\AtBeginDocument{%
\ifdefined\contentsname
  \renewcommand*\contentsname{Table of contents}
\else
  \newcommand\contentsname{Table of contents}
\fi
\ifdefined\listfigurename
  \renewcommand*\listfigurename{List of Figures}
\else
  \newcommand\listfigurename{List of Figures}
\fi
\ifdefined\listtablename
  \renewcommand*\listtablename{List of Tables}
\else
  \newcommand\listtablename{List of Tables}
\fi
\ifdefined\figurename
  \renewcommand*\figurename{Figure}
\else
  \newcommand\figurename{Figure}
\fi
\ifdefined\tablename
  \renewcommand*\tablename{Table}
\else
  \newcommand\tablename{Table}
\fi
}
\@ifpackageloaded{float}{}{\usepackage{float}}
\floatstyle{ruled}
\@ifundefined{c@chapter}{\newfloat{codelisting}{h}{lop}}{\newfloat{codelisting}{h}{lop}[chapter]}
\floatname{codelisting}{Listing}
\newcommand*\listoflistings{\listof{codelisting}{List of Listings}}
\makeatother
\makeatletter
\makeatother
\makeatletter
\@ifpackageloaded{caption}{}{\usepackage{caption}}
\@ifpackageloaded{subcaption}{}{\usepackage{subcaption}}
\makeatother

\usepackage{bookmark}
\IfFileExists{xurl.sty}{\usepackage{xurl}}{} % add URL line breaks if available
\urlstyle{same}
\hypersetup{
  pdftitle={Chapter 3: Risk Models},
  hidelinks,
  pdfcreator={LaTeX via pandoc}}


\title{\texorpdfstring{Chapter 3: Risk Models}{Chapter 3: Risk Models}}
\author{}
\date{}

\begin{document}
\frame{\titlepage}


\section{Compound Negative Binomial
Distributions}\label{compound-negative-binomial}

\begin{frame}{From Compound Poisson to Compound Negative Binomial}
\protect\phantomsection\label{from-compound-poisson-to-compound-negative-binomial}
In the collective risk model,

\[
S=\sum_{i=1}^{N}X_i.
\]

The distribution of \(S\) depends on both:

\begin{itemize}
\tightlist
\item
  the claim number distribution \(N\);
\item
  the claim severity distribution \(X\).
\end{itemize}

So far, we have considered the compound Poisson model.
\end{frame}

\begin{frame}{Why Do We Need Another Frequency Distribution?}
\protect\phantomsection\label{why-do-we-need-another-frequency-distribution}
For a Poisson random variable,

\[
E[N]=\operatorname{Var}(N)=\lambda.
\]

However, claim-count data may exhibit \textbf{overdispersion}:

\[
\operatorname{Var}(N)>E[N].
\]

In such cases, the Poisson distribution may be too restrictive.

The negative binomial distribution provides an alternative.
\end{frame}

\begin{frame}{The Negative Binomial Distribution --- Type 2}
\protect\phantomsection\label{the-negative-binomial-distribution-type-2}
Under the IFoA Type 2 convention,

\[
N\sim NB(k,p),
\]

where \(N\) is the number of \textbf{failures before the \(k\)th
success}.

The parameters satisfy

\[
k>0,\qquad 0<p<1,\qquad q=1-p.
\]

The support is

\[
N=0,1,2,\ldots
\]
\end{frame}

\begin{frame}{What Does the Negative Binomial Count?}
\protect\phantomsection\label{what-does-the-negative-binomial-count}
Imagine a sequence of independent Bernoulli trials.

We continue until the \(k\)th success occurs.

The random variable \(N\) counts the number of failures observed before
that success.

For example, when \(k=3\),

\[
F,F,S,F,F,S,F,S
\]

gives \(N=5\).
\end{frame}

\begin{frame}{Probability Function}
\protect\phantomsection\label{probability-function}
For \(n=0,1,2,\ldots\),

\[
\Pr(N=n)
=
\frac{\Gamma(k+n)}
{\Gamma(n+1)\Gamma(k)}
p^kq^n.
\]

When \(k\) is an integer,

\[
\Pr(N=n)
=
\binom{k+n-1}{n}p^kq^n.
\]

This is the \textbf{Negative Binomial Type 2} parameterisation.
\end{frame}

\begin{frame}{Properties of the Negative Binomial}
\protect\phantomsection\label{properties-of-the-negative-binomial}
The mean is

\[
E[N]=\frac{kq}{p}.
\]

The variance is

\[
\operatorname{Var}(N)=\frac{kq}{p^2}.
\]

Since \(0<p<1\),

\[
\operatorname{Var}(N)>E[N].
\]
\end{frame}

\begin{frame}{Probability Generating Function}
\protect\phantomsection\label{probability-generating-function}
The probability generating function is

\[
G_N(s)
=
E[s^N]
=
\left(\frac{p}{1-qs}\right)^k.
\]

The moment generating function is obtained by setting

\[
s=e^t.
\]

Hence,

\[
M_N(t)
=
\left(\frac{p}{1-qe^t}\right)^k.
\]
\end{frame}

\begin{frame}{Poisson versus Negative Binomial}
\protect\phantomsection\label{poisson-versus-negative-binomial}
{\def\LTcaptype{none} % do not increment counter
\begin{longtable}[]{@{}
  >{\raggedright\arraybackslash}p{(\linewidth - 4\tabcolsep) * \real{0.3333}}
  >{\raggedright\arraybackslash}p{(\linewidth - 4\tabcolsep) * \real{0.3333}}
  >{\raggedright\arraybackslash}p{(\linewidth - 4\tabcolsep) * \real{0.3333}}@{}}
\toprule\noalign{}
\begin{minipage}[b]{\linewidth}\raggedright
\end{minipage} & \begin{minipage}[b]{\linewidth}\raggedright
Poisson
\end{minipage} & \begin{minipage}[b]{\linewidth}\raggedright
Negative Binomial
\end{minipage} \\
\midrule\noalign{}
\endhead
Mean & \(\lambda\) & \(kq/p\) \\
Variance & \(\lambda\) & \(kq/p^2\) \\
Dispersion & \(\operatorname{Var}(N)=E[N]\) &
\(\operatorname{Var}(N)>E[N]\) \\
\bottomrule\noalign{}
\end{longtable}
}

The negative binomial allows \textbf{greater variability in claim
counts}.
\end{frame}

\begin{frame}{Why Use the Negative Binomial for Claim Counts?}
\protect\phantomsection\label{why-use-the-negative-binomial-for-claim-counts}
Insurance portfolios contain risks with different claim frequencies.

Consequently, observed claim counts may be more variable than allowed by
a Poisson model.

The negative binomial provides a flexible frequency model with

\[
\operatorname{Var}(N)>E[N].
\]

This makes it useful for modelling \textbf{overdispersed claim counts}.
\end{frame}

\begin{frame}{From Frequency to Aggregate Loss}
\protect\phantomsection\label{from-frequency-to-aggregate-loss}
The collective risk model is

\[
S=\sum_{i=1}^{N}X_i.
\]

Now let the claim number distribution be

\[
N\sim NB(k,p).
\]

The aggregate loss \(S\) is therefore a \textbf{compound negative
binomial random variable}.
\end{frame}

\begin{frame}{The Compound Negative Binomial Distribution}
\protect\phantomsection\label{the-compound-negative-binomial-distribution}
Let

\[
N\sim NB(k,p)
\]

and let

\[
X_1,X_2,\ldots\sim F_X
\]

be independent and identically distributed claim amounts, independent of
\(N\).

Then

\[
S=\sum_{i=1}^{N}X_i
\]

has a \textbf{compound negative binomial distribution}.
\end{frame}

\begin{frame}{Notation}
\protect\phantomsection\label{notation}
We denote the compound negative binomial distribution by

\[
S\sim\mathcal{CNB}(k,p,F_X).
\]

Here,

\begin{itemize}
\tightlist
\item
  \(k,p\) determine the claim frequency distribution;
\item
  \(F_X\) determines the claim severity distribution;
\item
  \(S\) represents aggregate claims.
\end{itemize}

The model separates \textbf{frequency} from \textbf{severity}.
\end{frame}

\begin{frame}{Mean of the Compound Negative Binomial}
\protect\phantomsection\label{mean-of-the-compound-negative-binomial}
For a compound distribution,

\[
E[S]=E[N]E[X].
\]

Let \(m_1=E[X].\)

Since

\[
E[N]=\frac{kq}{p},
\]

we obtain

\[
\boxed{E[S]=\frac{kq}{p}m_1}.
\]
\end{frame}

\begin{frame}{Variance of the Compound Negative Binomial}
\protect\phantomsection\label{variance-of-the-compound-negative-binomial}
For a compound distribution,

\[
\operatorname{Var}(S)
=
E[N]\operatorname{Var}(X)
+
\operatorname{Var}(N)(E[X])^2.
\]

Let

\[
m_1=E[X],
\qquad
m_2=E[X^2].
\]

Then

\[
\boxed{
\operatorname{Var}(S)
=
\frac{kq}{p^2}
\left(pm_2+qm_1^2\right)
}.
\]
\end{frame}

\begin{frame}{Interpreting the Variance}
\protect\phantomsection\label{interpreting-the-variance}
The aggregate loss variance reflects two sources of uncertainty:

\[
\operatorname{Var}(S)
=
\underbrace{E[N]\operatorname{Var}(X)}
_{\text{severity variation}}
+
\underbrace{\operatorname{Var}(N)(E[X])^2}
_{\text{frequency variation}}.
\]

Thus, both claim frequency and claim severity contribute to aggregate
risk.
\end{frame}

\begin{frame}{MGF of the Compound Distribution}
\protect\phantomsection\label{mgf-of-the-compound-distribution}
For a compound distribution,

\[
M_S(t)=G_N(M_X(t)).
\]

For the negative binomial,

\[
G_N(s)
=
\left(\frac{p}{1-qs}\right)^k.
\]

Therefore,

\[
\boxed{
M_S(t)
=
\left(
\frac{p}{1-qM_X(t)}
\right)^k
}.
\]
\end{frame}

\begin{frame}{Worked Example}
\protect\phantomsection\label{worked-example}
Suppose \(N\sim NB(k,p),\) with \(k=4,\qquad p=0.5.\)

Claim amounts have \(m_1=E[X]=1000,\qquad m_2=E[X^2]=1{,}500{,}000.\)

Calculate:

\begin{itemize}
\tightlist
\item
  the expected number of claims;
\item
  the variance of the number of claims;
\item
  the expected aggregate loss;
\item
  the variance of aggregate loss.
\end{itemize}
\end{frame}

\begin{frame}{Worked Example: Claim Number}
\protect\phantomsection\label{worked-example-claim-number}
Given \(k=4\), \(p=0.5\), and \(q=0.5\),

\[
E[N]=\frac{kq}{p}
=\frac{4(0.5)}{0.5}=4.
\]

\[
\operatorname{Var}(N)
=\frac{kq}{p^2}
=\frac{4(0.5)}{0.5^2}=8.
\]
\end{frame}

\begin{frame}{Worked Example: Aggregate Loss}
\protect\phantomsection\label{worked-example-aggregate-loss}
Given \(m_1=1{,}000\) and \(m_2=1{,}500{,}000\),

\[
E[S]=\frac{kq}{p}m_1
=4(1{,}000)=4{,}000.
\]

\[
\operatorname{Var}(S)
=\frac{kq}{p^2}(pm_2+qm_1^2)
=8[0.5(1{,}500{,}000)+0.5(1{,}000)^2]
=10{,}000{,}000.
\]
\end{frame}

\begin{frame}{R Simulation}
\protect\phantomsection\label{r-simulation}
We can simulate the collective risk model in two stages.

First, simulate the number of claims:

\[
N\sim NB(k,p).
\]

Then simulate the claim amounts

\[
X_1,\ldots,X_N\sim F_X.
\]

Finally, calculate

\[
S=\sum_{i=1}^{N}X_i.
\]
\end{frame}

\begin{frame}[fragile]{R Simulation: Setting Up the Model}
\protect\phantomsection\label{r-simulation-setting-up-the-model}
We simulate \(n_{\text{sim}}=5{,}000\) insurance years.

\begin{Shaded}
\begin{Highlighting}[]
\FunctionTok{set.seed}\NormalTok{(}\DecValTok{123}\NormalTok{)}

\NormalTok{k }\OtherTok{\textless{}{-}} \DecValTok{4}
\NormalTok{p }\OtherTok{\textless{}{-}} \FloatTok{0.5}
\NormalTok{q }\OtherTok{\textless{}{-}} \DecValTok{1} \SpecialCharTok{{-}}\NormalTok{ p}

\NormalTok{shape }\OtherTok{\textless{}{-}} \DecValTok{2}
\NormalTok{rate }\OtherTok{\textless{}{-}} \FloatTok{0.002}
\NormalTok{n\_sim }\OtherTok{\textless{}{-}} \DecValTok{5000}
\end{Highlighting}
\end{Shaded}

Here, \texttt{k} and \texttt{p} define the negative binomial frequency,
while \texttt{shape} and \texttt{rate} define the Gamma severity
distribution.
\end{frame}

\begin{frame}[fragile]{R Simulation: Claim Numbers}
\protect\phantomsection\label{r-simulation-claim-numbers}
First, simulate the number of claims in each year.

\begin{Shaded}
\begin{Highlighting}[]
\NormalTok{N }\OtherTok{\textless{}{-}} \FunctionTok{rnbinom}\NormalTok{(}
\NormalTok{  n\_sim,}
  \AttributeTok{size =}\NormalTok{ k,}
  \AttributeTok{prob =}\NormalTok{ p}
\NormalTok{)}
\end{Highlighting}
\end{Shaded}

Each element of \texttt{N} represents the number of claims in one
simulated year.

Thus,

\[
N\sim NB(k,p).
\]

For these parameters,

\[
E[N]=4,\qquad \operatorname{Var}(N)=8.
\]
\end{frame}

\begin{frame}[fragile]{R Simulation: Claim Amounts}
\protect\phantomsection\label{r-simulation-claim-amounts}
For each simulated year, generate the individual claim amounts.

\begin{Shaded}
\begin{Highlighting}[]
\NormalTok{S }\OtherTok{\textless{}{-}} \FunctionTok{sapply}\NormalTok{(N, }\ControlFlowTok{function}\NormalTok{(n) \{}
  \ControlFlowTok{if}\NormalTok{ (n }\SpecialCharTok{==} \DecValTok{0}\NormalTok{) }\FunctionTok{return}\NormalTok{(}\DecValTok{0}\NormalTok{)}
  \FunctionTok{sum}\NormalTok{(}\FunctionTok{rgamma}\NormalTok{(}
\NormalTok{    n, }\AttributeTok{shape =}\NormalTok{ shape, }\AttributeTok{rate =}\NormalTok{ rate}
\NormalTok{  ))}
\NormalTok{\})}
\end{Highlighting}
\end{Shaded}

For a year with \(N=n\) claims, this code generates \(n\) Gamma claim
amounts and adds them together.

Thus,

\[
S=\sum_{i=1}^{N}X_i.
\]
\end{frame}

\begin{frame}[fragile]{Claim Number: Simulation Result}
\protect\phantomsection\label{claim-number-simulation-result}
\begin{Shaded}
\begin{Highlighting}[]
\FunctionTok{hist}\NormalTok{(}
\NormalTok{  N,}
  \AttributeTok{breaks =} \FunctionTok{seq}\NormalTok{(}\SpecialCharTok{{-}}\FloatTok{0.5}\NormalTok{, }\FunctionTok{max}\NormalTok{(N) }\SpecialCharTok{+} \FloatTok{0.5}\NormalTok{, }\DecValTok{1}\NormalTok{),}
  \AttributeTok{main =} \StringTok{"Claim Number: Negative Binomial"}\NormalTok{,}
  \AttributeTok{xlab =} \StringTok{"Number of claims, N"}
\NormalTok{)}
\end{Highlighting}
\end{Shaded}

The histogram shows the simulated distribution of the annual claim
count.

Each bar represents the number of simulated years having a particular
value of \(N\).

The distribution includes the possibility of \textbf{zero claims}.
\end{frame}

\begin{frame}[fragile]{Aggregate Loss: Simulation Result}
\protect\phantomsection\label{aggregate-loss-simulation-result}
\begin{Shaded}
\begin{Highlighting}[]
\FunctionTok{hist}\NormalTok{(}
\NormalTok{  S,}
  \AttributeTok{breaks =} \DecValTok{50}\NormalTok{,}
  \AttributeTok{main =} \StringTok{"Aggregate Loss: Compound Negative Binomial"}\NormalTok{,}
  \AttributeTok{xlab =} \StringTok{"Aggregate loss, S"}
\NormalTok{)}
\end{Highlighting}
\end{Shaded}

The histogram shows the simulated distribution of annual aggregate
losses.

Each value of \texttt{S} represents one simulated insurance year.

The right-skewed shape arises from variation in both \textbf{claim
frequency} and \textbf{claim severity}.
\end{frame}

\begin{frame}[fragile]{Simulation: From Frequency to Aggregate Loss}
\protect\phantomsection\label{simulation-from-frequency-to-aggregate-loss}
The simulation follows the compound model in three stages:

\[
N\sim NB(k,p)
\]

\[
X_1,\ldots,X_N\sim F_X
\]

\[
S=\sum_{i=1}^{N}X_i.
\]

Thus, each element of \texttt{S} represents \textbf{one simulated
aggregate loss}.
\end{frame}

\begin{frame}[fragile]{Simulation versus Theory}
\protect\phantomsection\label{simulation-versus-theory}
We can calculate the theoretical severity moments.

\begin{Shaded}
\begin{Highlighting}[]
\NormalTok{m1 }\OtherTok{\textless{}{-}}\NormalTok{ shape }\SpecialCharTok{/}\NormalTok{ rate}
\NormalTok{m2 }\OtherTok{\textless{}{-}}\NormalTok{ shape }\SpecialCharTok{*}\NormalTok{ (shape }\SpecialCharTok{+} \DecValTok{1}\NormalTok{) }\SpecialCharTok{/}\NormalTok{ rate}\SpecialCharTok{\^{}}\DecValTok{2}

\NormalTok{theory\_mean }\OtherTok{\textless{}{-}}\NormalTok{ k }\SpecialCharTok{*}\NormalTok{ q }\SpecialCharTok{/}\NormalTok{ p }\SpecialCharTok{*}\NormalTok{ m1}
\NormalTok{theory\_var }\OtherTok{\textless{}{-}}\NormalTok{ k }\SpecialCharTok{*}\NormalTok{ q }\SpecialCharTok{/}\NormalTok{ p}\SpecialCharTok{\^{}}\DecValTok{2} \SpecialCharTok{*}
\NormalTok{  (p }\SpecialCharTok{*}\NormalTok{ m2 }\SpecialCharTok{+}\NormalTok{ q }\SpecialCharTok{*}\NormalTok{ m1}\SpecialCharTok{\^{}}\DecValTok{2}\NormalTok{)}
\end{Highlighting}
\end{Shaded}

Here, \texttt{m1} and \texttt{m2} are the first and second moments of
the Gamma severity.

These are substituted into the compound negative binomial formulas.
\end{frame}

\begin{frame}[fragile]{Simulation versus Theory}
\protect\phantomsection\label{simulation-versus-theory-1}
Compare the simulated and theoretical results.

\begin{Shaded}
\begin{Highlighting}[]
\FunctionTok{c}\NormalTok{(}
  \AttributeTok{Simulated\_Mean =} \FunctionTok{mean}\NormalTok{(S),}
  \AttributeTok{Theoretical\_Mean =}\NormalTok{ theory\_mean,}
  \AttributeTok{Simulated\_Var =} \FunctionTok{var}\NormalTok{(S),}
  \AttributeTok{Theoretical\_Var =}\NormalTok{ theory\_var}
\NormalTok{)}
\end{Highlighting}
\end{Shaded}

The simulated values should be close to the theoretical values.

With many simulated years, the differences become smaller.

This provides a numerical check of the compound negative binomial
formulas.
\end{frame}

\begin{frame}{Theoretical Aggregate Loss}
\protect\phantomsection\label{theoretical-aggregate-loss}
For the Gamma severity distribution,

\[
E[X]=\frac{\text{shape}}{\text{rate}}=1000
\]

and

\[
E[X^2]
=
\frac{\text{shape}(\text{shape}+1)}
{\text{rate}^2}
=1{,}500{,}000.
\]

Therefore,

\[
E[S]
=
\frac{kq}{p}m_1
=
4(1000)
=
4000.
\]

The simulation should produce an aggregate mean close to \(4000\).
\end{frame}

\begin{frame}{Graph: Aggregate Loss}
\protect\phantomsection\label{graph-aggregate-loss}
\pandocbounded{\includegraphics[keepaspectratio]{04_Slides_Risk_Models_Part2_files/figure-beamer/unnamed-chunk-1-1.pdf}}
\end{frame}

\begin{frame}[fragile]{Compare Simulation with Theory}
\protect\phantomsection\label{compare-simulation-with-theory}
And similarly, this can be a \textbf{self-contained chunk} that
calculates the theoretical moments and compares them with the
simulation:

\begin{verbatim}
   Measure  Simulation Theory
1     Mean    3926.812  4e+03
2 Variance 9437849.182  1e+07
\end{verbatim}
\end{frame}

\begin{frame}{Interpreting the Simulation}
\protect\phantomsection\label{interpreting-the-simulation}
The simulated and theoretical results are reasonably close:

{\def\LTcaptype{none} % do not increment counter
\begin{longtable}[]{@{}lrr@{}}
\toprule\noalign{}
Measure & Simulation & Theory \\
\midrule\noalign{}
\endhead
Mean & 3,926.8 & 4,000 \\
Variance & 9,437,849.2 & 10,000,000 \\
\bottomrule\noalign{}
\end{longtable}
}

\begin{itemize}
\tightlist
\item
  The simulated mean is close to the theoretical mean of \(4,000\).
\item
  The simulated variance is also reasonably close to the theoretical
  value.
\item
  The differences arise because only \(5,000\) insurance years were
  simulated.
\item
  Increasing the number of simulations would generally make the
  simulated results closer to the theoretical values.
\item
  This provides a numerical check of the compound negative binomial
  formulas.
\end{itemize}

The key idea is that \textbf{simulation reproduces the theoretical
behaviour of the model approximately}.
\end{frame}

\begin{frame}{When Might We Use This Model?}
\protect\phantomsection\label{when-might-we-use-this-model}
The compound negative binomial model is useful when:

\begin{itemize}
\tightlist
\item
  claim counts are overdispersed;
\item
  \(\operatorname{Var}(N)>E[N]\);
\item
  the Poisson model is too restrictive;
\item
  a flexible frequency model is required;
\item
  aggregate claim amounts are the quantity of interest.
\end{itemize}
\end{frame}

\begin{frame}{Limitations}
\protect\phantomsection\label{limitations}
The negative binomial distribution is not automatically appropriate
whenever

\[
\operatorname{Var}(N)>E[N].
\]

The choice of frequency distribution should be supported by:

\begin{itemize}
\tightlist
\item
  the characteristics of the data;
\item
  goodness-of-fit analysis;
\item
  the modelling assumptions;
\item
  the intended application.
\end{itemize}
\end{frame}

\begin{frame}{Summary}
\protect\phantomsection\label{summary}
The key progression is

\[
\text{Poisson}
\rightarrow
\text{Overdispersion}
\rightarrow
\text{Negative Binomial}
\rightarrow
\text{Compound Negative Binomial}.
\]

For \(N\sim NB(k,p)\),

\[
E[N]=\frac{kq}{p},
\qquad
\operatorname{Var}(N)=\frac{kq}{p^2}.
\]
\end{frame}

\begin{frame}{Summary: Aggregate Loss}
\protect\phantomsection\label{summary-aggregate-loss}
For \(S=\sum_{i=1}^{N}X_i,\)

the compound negative binomial model gives

\[
E[S]=\frac{kq}{p}m_1,
\]

and

\[
\operatorname{Var}(S)
=
\frac{kq}{p^2}
\left(pm_2+qm_1^2\right).
\]

Its MGF is

\[
M_S(t)
=
\left(
\frac{p}{1-qM_X(t)}
\right)^k.
\]
\end{frame}

\section{Compound Binomial Distributions}\label{compound-binomial}

\begin{frame}{From Individual Policies to Aggregate Loss}
\protect\phantomsection\label{from-individual-policies-to-aggregate-loss}
Consider a portfolio of \(n\) independent and identical policies.

Each policy can give rise to \textbf{at most one claim} during the
period.

Let \(p\) be the probability that a policy produces a claim.

The number of claims in the portfolio is therefore

\[
N \sim \operatorname{Binomial}(n,p).
\]
\end{frame}

\begin{frame}{Why the Binomial Distribution?}
\protect\phantomsection\label{why-the-binomial-distribution}
The binomial model is natural when:

\begin{itemize}
\tightlist
\item
  there are a fixed number \(n\) of potential claim events;
\item
  each policy produces either \textbf{zero or one claim};
\item
  each policy has the same claim probability \(p\);
\item
  claim events are independent.
\end{itemize}

Thus, the number of claims has an upper limit:

\[
0 \leq N \leq n.
\]
\end{frame}

\begin{frame}{The Compound Binomial Model}
\protect\phantomsection\label{the-compound-binomial-model}
Let \(X_i\) denote the amount of the \(i\)th claim.

Assume that \(X_1,X_2,\ldots\) are i.i.d. and independent of \(N\).

Then the aggregate claim amount is

\[
S=\sum_{i=1}^{N}X_i.
\]

We call \(S\) a \textbf{compound binomial random variable}.

The distribution is denoted by

\[
S\sim\mathcal{CB}(n,p,F_X).
\]
\end{frame}

\begin{frame}{Aggregate Loss Model}
\protect\phantomsection\label{aggregate-loss-model}
The collective risk model is

\[
S=\sum_{i=1}^{N}X_i,
\]

where

\[
N\sim\operatorname{Binomial}(n,p).
\]

Here, \(N\) determines \textbf{how many claims occur}, while \(X_i\)
determines \textbf{how much each claim costs}.

Thus, the aggregate loss combines a \textbf{frequency model} and a
\textbf{severity model}.
\end{frame}

\begin{frame}{Expected Aggregate Loss}
\protect\phantomsection\label{expected-aggregate-loss}
For a random sum,

\[
S=\sum_{i=1}^{N}X_i,
\]

the expected aggregate loss is

\[
\boxed{E[S]=E[N]E[X]}.
\]

This follows from the standard result for a random sum when \(N\) is
independent of the claim amounts.
\end{frame}

\begin{frame}{Binomial Frequency Model}
\protect\phantomsection\label{binomial-frequency-model}
For the binomial frequency model,

\[
N\sim\operatorname{Binomial}(n,p),
\]

so

\[
E[N]=np.
\]

Let

\[
m_1=E[X].
\]

Substituting into the random-sum formula,

\[
\begin{aligned}
E[S]
&=E[N]E[X]\\
&=(np)m_1.
\end{aligned}
\]

Therefore,

\[
\boxed{E[S]=npm_1}.
\]
\end{frame}

\begin{frame}{Interpretation of Expected Aggregate Loss}
\protect\phantomsection\label{interpretation-of-expected-aggregate-loss}
The formula

\[
E[S]=npm_1
\]

can be written as

\[
\boxed{
E[S]
=
\underbrace{np}_{\text{expected number of claims}}
\underbrace{m_1}_{\text{expected claim amount}}
}.
\]

Thus, expected aggregate loss is determined by:

\begin{itemize}
\tightlist
\item
  the expected number of claims, \(np\);
\item
  the expected amount per claim, \(m_1\).
\end{itemize}

The binomial frequency model therefore gives a simple interpretation of
the expected aggregate loss.
\end{frame}

\begin{frame}{Variance of Aggregate Loss}
\protect\phantomsection\label{variance-of-aggregate-loss}
For a random sum,

\[
\operatorname{Var}(S)
=
E[N]\operatorname{Var}(X)
+
\operatorname{Var}(N)\{E[X]\}^2.
\]

For

\[
N\sim\operatorname{Binomial}(n,p),
\]

we have

\[
E[N]=np,
\qquad
\operatorname{Var}(N)=np(1-p).
\]

Let

\[
m_1=E[X],
\qquad
m_2=E[X^2].
\]
\end{frame}

\begin{frame}{Substituting the Binomial Moments}
\protect\phantomsection\label{substituting-the-binomial-moments}
Since

\[
\operatorname{Var}(X)=m_2-m_1^2,
\]

substitute into the random-sum formula:

\[
\begin{aligned}
\operatorname{Var}(S)
&=np(m_2-m_1^2)
+np(1-p)m_1^2.
\end{aligned}
\]

Expanding,

\[
\begin{aligned}
\operatorname{Var}(S)
&=npm_2-npm_1^2
+np(1-p)m_1^2.
\end{aligned}
\]
\end{frame}

\begin{frame}{Simplifying the Variance}
\protect\phantomsection\label{simplifying-the-variance}
Collect the terms involving \(m_1^2\):

\[
\begin{aligned}
\operatorname{Var}(S)
&=npm_2
-npm_1^2
+np(1-p)m_1^2\\
&=npm_2-np^2m_1^2.
\end{aligned}
\]

Therefore,

\[
\boxed{
\operatorname{Var}(S)
=
npm_2-np^2m_1^2
}.
\]

The result combines \textbf{claim-size uncertainty} and
\textbf{frequency uncertainty} under the binomial model.
\end{frame}

\begin{frame}{Worked Example: Compound Binomial with Gamma Severity}
\protect\phantomsection\label{worked-example-compound-binomial-with-gamma-severity}
Suppose

\[
N\sim\operatorname{Binomial}(100,0.01),
\]

and, independently,

\[
X_i\overset{\text{i.i.d.}}{\sim}\operatorname{Gamma}(20,0.1),
\]

where the Gamma distribution uses the \textbf{shape-rate}
parameterisation.

Let

\[
S=\sum_{i=1}^{N}X_i.
\]
\end{frame}

\begin{frame}{Worked Example: Tasks}
\protect\phantomsection\label{worked-example-tasks}
For the compound binomial model above, find:

\begin{enumerate}
\tightlist
\item
  the MGF of \(S\);
\item
  \(E[S]\) using the MGF;
\item
  \(\operatorname{Var}(S)\) using the MGF;
\item
  \(P(S=0)\).
\end{enumerate}
\end{frame}

\begin{frame}{1. MGF of the Aggregate Loss}
\protect\phantomsection\label{mgf-of-the-aggregate-loss}
For

\[
X\sim\operatorname{Gamma}(20,0.1),
\]

the MGF is

\[
M_X(t)
=
\left(\frac{0.1}{0.1-t}\right)^{20}.
\]

For the compound binomial model,

\[
M_S(t)
=
\left[q+pM_X(t)\right]^n.
\]

Here, \(n=100,\qquad p=0.01,\qquad q=0.99.\)
\end{frame}

\begin{frame}{MGF of \(S\)}
\protect\phantomsection\label{mgf-of-s}
Substituting the values,

\[
\boxed{
M_S(t)=
\left[
0.99+
0.01\left(\frac{0.1}{0.1-t}\right)^{20}
\right]^{100}
}
\]

for \(t<0.1\).

This follows directly from the compound binomial MGF.
\end{frame}

\begin{frame}{2. Expected Aggregate Loss}
\protect\phantomsection\label{expected-aggregate-loss-1}
Use \(E[S]=M_S'(0).\)

Let

\[
A(t)=
0.99+
0.01\left(\frac{0.1}{0.1-t}\right)^{20}.
\]

Then

\[
M_S(t)=A(t)^{100}.
\]

Hence, \(M_S'(t)=100A(t)^{99}A'(t).\)

At \(t=0\), \qquad \(A(0)=1, \qquad A'(0)=0.01M_X'(0).\)
\end{frame}

\begin{frame}{Expected Aggregate Loss}
\protect\phantomsection\label{expected-aggregate-loss-2}
For the Gamma distribution,

\[
M_X'(0)=E[X]
=\frac{20}{0.1}=200.
\]

Therefore,

\[
\begin{aligned}
E[S]
&=M_S'(0)\\
&=100(0.01)(200).
\end{aligned}
\]

Thus,

\[
\boxed{E[S]=200}.
\]
\end{frame}

\begin{frame}{3. Variance Using the MGF}
\protect\phantomsection\label{variance-using-the-mgf}
Use

\[
\operatorname{Var}(S)
=
M_S''(0)-[M_S'(0)]^2.
\]

From the previous result,

\[
M_S'(0)=200.
\]

We therefore need

\[
M_S''(0).
\]
\end{frame}

\begin{frame}{Second Derivative}
\protect\phantomsection\label{second-derivative}
For

\[
M_S(t)=A(t)^{100},
\]

the second derivative is

\[
M_S''(t)
=
100(99)A(t)^{98}[A'(t)]^2
+
100A(t)^{99}A''(t).
\]

At \(t=0\),

\[
M_S''(0)
=
100(99)[A'(0)]^2
+
100A''(0).
\]
\end{frame}

\begin{frame}{Evaluate the Components}
\protect\phantomsection\label{evaluate-the-components}
For the Gamma severity,

\[
E[X^2]
=
\operatorname{Var}(X)+[E(X)]^2.
\]

Since

\[
\operatorname{Var}(X)=\frac{20}{0.1^2}=2{,}000,
\]

we obtain

\[
E[X^2]=2{,}000+200^2=42{,}000.
\]

Therefore, \(A'(0)=0.01(200), \qquad A''(0)=0.01(42{,}000).\)
\end{frame}

\begin{frame}{Variance of \(S\)}
\protect\phantomsection\label{variance-of-s}
Thus,

\[
M_S''(0)
=
100(99)(0.01\times200)^2
+
100(0.01\times42{,}000).
\]

Hence,

\[
M_S''(0)=81{,}600.
\]

Finally,

\[
\begin{aligned}
\operatorname{Var}(S)
&=81{,}600-200^2\\
&=\boxed{41{,}600}.
\end{aligned}
\]
\end{frame}

\begin{frame}{4. Probability of Zero Aggregate Loss}
\protect\phantomsection\label{probability-of-zero-aggregate-loss}
Because claim amounts are positive,

\[
S=0
\quad\Longleftrightarrow\quad
N=0.
\]

Therefore,

\[
P(S=0)=P(N=0).
\]

For

\[
N\sim\operatorname{Binomial}(100,0.01),
\]

we have

\[
P(N=0)
=
(1-0.01)^{100}.
\]
\end{frame}

\begin{frame}{Evaluating \(P(S=0)\)}
\protect\phantomsection\label{evaluating-ps0}
Thus,

\[
P(S=0)
=
(0.99)^{100}.
\]

Numerically,

\[
\boxed{P(S=0)\approx0.3660}.
\]

So there is approximately a \textbf{36.6\% probability of no aggregate
claims}.
\end{frame}

\begin{frame}{An Alternative MGF View}
\protect\phantomsection\label{an-alternative-mgf-view}
Recall that \(M_S(t)=E[e^{tS}].\)

As \(t\to-\infty\),

\[
e^{tS}\to
\begin{cases}
1, & S=0,\\
0, & S>0.
\end{cases}
\]

Therefore,

\[
P(S=0)
=
\lim_{t\to-\infty}M_S(t).
\]

For this model,
\(\qquad \lim_{t\to-\infty}M_S(t)=(0.99)^{100} \approx0.3660.\)

This provides another useful interpretation of the MGF.
\end{frame}

\begin{frame}{Key Lessons}
\protect\phantomsection\label{key-lessons}
The compound binomial MGF is obtained by substituting the severity MGF
into

\[
M_S(t)=\left[q+pM_X(t)\right]^n.
\]

For this example,
\(\boxed{M_S(t)=\left[0.99+0.01\left(\frac{0.1}{0.1-t}\right)^{20}\right]^{100}}.\)

The MGF can be used to obtain moments:

\[
E[S]=200,
\qquad
\operatorname{Var}(S)=41{,}600.
\]

The model has a positive probability of zero aggregate loss:
\(P(S=0)=(0.99)^{100}\approx0.3660.\)

\textbf{Main idea:} the compound-binomial MGF provides a single
framework for obtaining both \textbf{aggregate-loss moments} and other
distributional information.
\end{frame}

\begin{frame}{R Simulation}
\protect\phantomsection\label{r-simulation-1}
We can simulate the collective risk model in two stages.

\begin{enumerate}
\item
  \textbf{Simulate the number of claims} \(N\).
\item
  \textbf{Simulate the claim amounts} \(X_1,\ldots,X_N\) and calculate
  the aggregate loss

  \[
  S=\sum_{i=1}^{N}X_i.
  \]
\end{enumerate}

For this example,

\[
N\sim\operatorname{Binomial}(100,0.01)
\]

and

\[
X_i\overset{\text{i.i.d.}}{\sim}\operatorname{Gamma}(20,0.1).
\]
\end{frame}

\begin{frame}[fragile]{Stage 1: Simulate the Number of Claims}
\protect\phantomsection\label{stage-1-simulate-the-number-of-claims}
For each simulated year, generate

\[
N\sim\operatorname{Binomial}(100,0.01).
\]

Each element of \texttt{N} represents the \textbf{number of claims in
one simulated year}.

For example,

\begin{verbatim}
[1] 0 2 1 2 3 0
\end{verbatim}

The simulation produces \(B=100{,}000\) independent insurance years.
\end{frame}

\begin{frame}[fragile]{Stage 2: Simulate Aggregate Loss}
\protect\phantomsection\label{stage-2-simulate-aggregate-loss}
For each simulated year:

\begin{itemize}
\tightlist
\item
  If \(N=0\), then \(S=0\).
\item
  If \(N=n>0\), simulate \(n\) claim amounts.
\item
  Add the claim amounts to obtain \(S\).
\end{itemize}

Thus, each element of \texttt{S} represents one simulated value of

\[
S=\sum_{i=1}^{N}X_i.
\]
\end{frame}

\begin{frame}[fragile]{R Simulation}
\protect\phantomsection\label{r-simulation-2}
We can simulate the aggregate loss directly in R.

\begin{Shaded}
\begin{Highlighting}[]
\CommentTok{\# Simulate aggregate losses}
\FunctionTok{set.seed}\NormalTok{(}\DecValTok{123}\NormalTok{)}
\NormalTok{B }\OtherTok{\textless{}{-}} \DecValTok{100000}
\NormalTok{N }\OtherTok{\textless{}{-}} \FunctionTok{rbinom}\NormalTok{(B, }\AttributeTok{size =} \DecValTok{100}\NormalTok{, }\AttributeTok{prob =} \FloatTok{0.01}\NormalTok{)}
\NormalTok{S }\OtherTok{\textless{}{-}} \FunctionTok{sapply}\NormalTok{(N, }\ControlFlowTok{function}\NormalTok{(n) \{}
  \FunctionTok{sum}\NormalTok{(}\FunctionTok{rgamma}\NormalTok{(n, }\AttributeTok{shape =} \DecValTok{20}\NormalTok{, }\AttributeTok{rate =} \FloatTok{0.1}\NormalTok{))}
\NormalTok{\})}

\CommentTok{\# Plot the simulated aggregate losses}
\FunctionTok{hist}\NormalTok{(S, }\AttributeTok{breaks =} \DecValTok{50}\NormalTok{,}
     \AttributeTok{main =} \StringTok{"Simulated Aggregate Loss"}\NormalTok{,}
     \AttributeTok{xlab =} \StringTok{"Aggregate loss, S"}\NormalTok{)}
\end{Highlighting}
\end{Shaded}
\end{frame}

\begin{frame}{R Simulation}
\protect\phantomsection\label{r-simulation-3}
We can simulate the aggregate loss directly in R.

\begin{center}
\pandocbounded{\includegraphics[keepaspectratio]{04_Slides_Risk_Models_Part2_files/figure-beamer/unnamed-chunk-6-1.pdf}}
\end{center}
\end{frame}

\begin{frame}[fragile]{Simulation Results}
\protect\phantomsection\label{simulation-results}
Use the simulated aggregate losses to estimate the moments.

\begin{verbatim}
[1] 199.89
\end{verbatim}

\begin{verbatim}
[1] 41515.39
\end{verbatim}

For \(100{,}000\) simulations, \(\widehat{E[S]}=199.89\) and
\(\widehat{\operatorname{Var}(S)}=41{,}515.39\).

Compare with the theoretical results:

\[
E[S]=200,
\qquad
\operatorname{Var}(S)=41{,}600.
\]

\textbf{The simulation closely reproduces the theoretical results.}
\end{frame}

\begin{frame}{From Model to Simulation}
\protect\phantomsection\label{from-model-to-simulation}
The simulation follows exactly the structure of the collective risk
model:

\[
\boxed{
N
\longrightarrow
X_1,\ldots,X_N
\longrightarrow
S
}
\] \textbf{Frequency:} \(N\sim\operatorname{Binomial}(100,0.01)\)

\textbf{Severity:} \(X_i\sim\operatorname{Gamma}(20,0.1)\)

\textbf{Aggregate loss:} \(S=\sum_{i=1}^{N}X_i\)

\begin{quote}
\textbf{Key idea:} simulation lets us generate many possible insurance
years and use the simulated aggregate losses to approximate quantities
such as \(E[S]\) and \(\operatorname{Var}(S)\).
\end{quote}
\end{frame}

\section{Properties of the Compound Binomial
Model}\label{properties-of-the-compound-binomial-model}

\begin{frame}{Comparing the Three Frequency Models}
\protect\phantomsection\label{comparing-the-three-frequency-models}
The three models differ mainly in their assumptions about the number of
claims \(N\).

{\def\LTcaptype{none} % do not increment counter
\begin{longtable}[]{@{}
  >{\raggedright\arraybackslash}p{(\linewidth - 4\tabcolsep) * \real{0.3562}}
  >{\raggedright\arraybackslash}p{(\linewidth - 4\tabcolsep) * \real{0.4521}}
  >{\raggedright\arraybackslash}p{(\linewidth - 4\tabcolsep) * \real{0.1918}}@{}}
\toprule\noalign{}
\begin{minipage}[b]{\linewidth}\raggedright
Model
\end{minipage} & \begin{minipage}[b]{\linewidth}\raggedright
Frequency distribution
\end{minipage} & \begin{minipage}[b]{\linewidth}\raggedright
Support of \(N\)
\end{minipage} \\
\midrule\noalign{}
\endhead
Compound Binomial & \(\operatorname{Binomial}(n,p)\) & \(0,\ldots,n\) \\
Compound Poisson & \(\operatorname{Poisson}(\lambda)\) &
\(0,1,2,\ldots\) \\
Compound Negative Binomial & \(\operatorname{NB}(k,p)\) &
\(0,1,2,\ldots\) \\
\bottomrule\noalign{}
\end{longtable}
}

The claim amounts \(X_i\) can have the same severity distribution in all
three models.
\end{frame}

\begin{frame}{Comparing Frequency Variability}
\protect\phantomsection\label{comparing-frequency-variability}
For the binomial model,

\[
E[N]=np,
\qquad
\operatorname{Var}(N)=np(1-p).
\]

For the Poisson model,

\[
E[N]=\lambda,
\qquad
\operatorname{Var}(N)=\lambda.
\]

For the negative binomial model,

\[
E[N]=\frac{kq}{p},
\qquad
\operatorname{Var}(N)=\frac{kq}{p^2}.
\]

Thus, the binomial model has \(\operatorname{Var}(N)<E[N],\)

while the Poisson and negative binomial models have variance equal to or
greater than their mean.
\end{frame}

\begin{frame}{Interpretation of Frequency Variability}
\protect\phantomsection\label{interpretation-of-frequency-variability}
The three models represent different patterns of claim-count
variability.

\textbf{Binomial:} fixed number of potential claim events and limited
frequency variability.

\textbf{Poisson:} no fixed upper limit, with frequency variance equal to
the mean.

\textbf{Negative binomial:} no fixed upper limit and greater frequency
variability than the Poisson model.

The choice of frequency model should reflect the underlying
claim-generation mechanism.
\end{frame}

\begin{frame}{A Fixed Upper Limit on Claims}
\protect\phantomsection\label{a-fixed-upper-limit-on-claims}
A defining feature of the compound binomial model is

\[
N\leq n.
\]

Consequently, the portfolio can contain at most \(n\) claims.

This is natural when there are \(n\) policies and each policy can
produce at most one claim.

In contrast, Poisson and negative binomial models have no finite upper
limit on \(N\).
\end{frame}

\begin{frame}{The Role of the Claim Probability}
\protect\phantomsection\label{the-role-of-the-claim-probability}
For

\[
N\sim\operatorname{Binomial}(n,p),
\]

the expected number of claims is

\[
E[N]=np.
\]

Increasing \(p\) increases the expected frequency.

However,\(\operatorname{Var}(N)=np(1-p),\) so frequency variability is
also affected by \(p\).

The variance is largest when \(p=0.5\).
\end{frame}

\begin{frame}{Aggregate Loss Depends on Frequency and Severity}
\protect\phantomsection\label{aggregate-loss-depends-on-frequency-and-severity}
For the compound binomial model,

\[
E[S]=npE[X].
\]

The expected aggregate loss therefore depends on:

\begin{itemize}
\tightlist
\item
  the number of potential claims \(n\);
\item
  the probability of a claim \(p\);
\item
  the expected claim amount \(E[X]\).
\end{itemize}

Similarly,

\[
\operatorname{Var}(S)
=
E[N]\operatorname{Var}(X)
+
\operatorname{Var}(N)\{E[X]\}^2.
\]

Thus, aggregate risk comes from both \textbf{claim severity} and
\textbf{claim frequency}.
\end{frame}

\begin{frame}{Binomial to Poisson as an Approximation}
\protect\phantomsection\label{binomial-to-poisson-as-an-approximation}
The binomial frequency model approaches a Poisson model when

\[
n\to\infty,
\qquad
p\to0,
\qquad
np\to\lambda.
\]

Then

\[
\operatorname{Binomial}(n,p)
\approx
\operatorname{Poisson}(\lambda).
\]

This provides a useful connection between the compound binomial and
compound Poisson models.

The Poisson model can therefore be viewed as an approximation when there
are many potential claim events with small individual claim
probabilities.
\end{frame}

\begin{frame}{Compound Binomial: Key Results}
\protect\phantomsection\label{compound-binomial-key-results}
For

\[
N\sim\operatorname{Binomial}(n,p),
\]

and

\[
S=\sum_{i=1}^{N}X_i,
\]

with \(N\) independent of the i.i.d. claim amounts \(X_i\),

\[
E[S]=np\,m_1,\qquad \operatorname{Var}(S)=npm_2-np^2m_1^2.
\]

The compound binomial model is particularly appropriate when there is a
\textbf{fixed number of potential claim events}, with each event
producing at most one claim.
\end{frame}

\begin{frame}{Three Compound Models}
\protect\phantomsection\label{three-compound-models}
The main distinction among the three special compound models is their
\textbf{frequency distribution}.

\[
\begin{array}{c|c|c}
\text{Model} & E[N] & \operatorname{Var}(N)\\
\hline
\text{Binomial} & np & np(1-p)\\
\text{Poisson} & \lambda & \lambda\\
\text{Negative Binomial} & \dfrac{kq}{p} & \dfrac{kq}{p^2}
\end{array}
\]

The frequency assumption determines the variability of the aggregate
loss.
\end{frame}

\section{Aggregate Claims with
Reinsurance}\label{aggregate-claims-with-reinsurance}

\begin{frame}{Next: Aggregate Claims with Reinsurance}
\protect\phantomsection\label{next-aggregate-claims-with-reinsurance}
We now extend the collective risk model to incorporate
\textbf{reinsurance}.

The next section considers aggregate claim distributions under:

\begin{itemize}
\tightlist
\item
  \textbf{Proportional reinsurance}
\item
  \textbf{Individual excess of loss reinsurance}
\end{itemize}

We will examine how reinsurance changes the \textbf{individual claim
amounts} and, consequently, the \textbf{aggregate claim distribution}.
\end{frame}

\begin{frame}{Proportional Reinsurance}
\protect\phantomsection\label{proportional-reinsurance}
Let \(S\) denote the total aggregate claims from a risk during a given
period.

Let \(S_I\) and \(S_R\) denote the aggregate claims paid by the
\textbf{insurer} and \textbf{reinsurer}, respectively.

Then

\[
\boxed{S=S_I+S_R}.
\]

Under proportional reinsurance, the insurer retains a fixed proportion
\(\alpha\) of every claim.
\end{frame}

\begin{frame}{Proportional Reinsurance}
\protect\phantomsection\label{proportional-reinsurance-1}
For each claim \(X_i\),

\[
Y_i=\alpha X_i,
\qquad
Z_i=(1-\alpha)X_i.
\]

Therefore,

\[
S_I=\sum_{i=1}^N\alpha X_i=\alpha S
\]

and

\[
S_R=\sum_{i=1}^N(1-\alpha)X_i=(1-\alpha)S.
\]
\end{frame}

\begin{frame}{Key Features}
\protect\phantomsection\label{key-features}
\begin{itemize}
\tightlist
\item
  Both the insurer and reinsurer contribute to \textbf{every claim}.
\item
  The insurer retains the same proportion \(\alpha\) of every claim.
\item
  The reinsurer pays the remaining proportion \(1-\alpha\).
\item
  Both parties have \textbf{unlimited liability} unless a separate cap
  is imposed.
\end{itemize}

Thus,

\[
\boxed{S_I=\alpha S,\qquad S_R=(1-\alpha)S}.
\]
\end{frame}

\begin{frame}{Worked Example}
\protect\phantomsection\label{worked-example-1}
Suppose aggregate claims follow a compound Poisson distribution with

\[
N\sim\operatorname{Poisson}(10)
\]

and individual claim amounts

\[
X\sim Pa(4,1).
\]

The insurer purchases proportional reinsurance with retention
\(\alpha=0.8.\)

Find:

\begin{enumerate}
\tightlist
\item
  the distributions of \(S_I\) and \(S_R\);
\item
  their means and variances;
\item
  compare \(\operatorname{Var}(S_I)+\operatorname{Var}(S_R)\) with
  \(\operatorname{Var}(S)\).
\end{enumerate}
\end{frame}

\begin{frame}{Solution: Set Up the Reinsurance Model}
\protect\phantomsection\label{solution-set-up-the-reinsurance-model}
The original aggregate claim is

\[
S=\sum_{i=1}^{N}X_i.
\]

Under proportional reinsurance with \(\alpha=0.8\),

\[
S_I=0.8S,
\qquad
S_R=0.2S.
\]

Equivalently,

\[
S_I=\sum_{i=1}^{N}0.8X_i,
\qquad
S_R=\sum_{i=1}^{N}0.2X_i.
\]
\end{frame}

\begin{frame}{Solution: Claim Amounts After Reinsurance}
\protect\phantomsection\label{solution-claim-amounts-after-reinsurance}
Define the insurer's payment per claim as

\[
Y_i=0.8X_i.
\]

Define the reinsurer's payment per claim as

\[
Z_i=0.2X_i.
\]

Then

\[
S_I=\sum_{i=1}^{N}Y_i,
\qquad
S_R=\sum_{i=1}^{N}Z_i.
\]

Both are therefore compound Poisson distributions.
\end{frame}

\begin{frame}{Solution: Distribution of the Insurer's Claim}
\protect\phantomsection\label{solution-distribution-of-the-insurers-claim}
Recall the scaling property:

\[
X\sim Pa(\alpha,\lambda)
\quad\Longrightarrow\quad
kX\sim Pa(\alpha,k\lambda).
\]

Since

\[
X_i\sim Pa(4,1)
\]

and

\[
Y_i=0.8X_i,
\]

we obtain

\[
\boxed{Y_i\sim Pa(4,0.8)}.
\]
\end{frame}

\begin{frame}{Solution: Distribution of the Reinsurer's Claim}
\protect\phantomsection\label{solution-distribution-of-the-reinsurers-claim}
Similarly,

\[
Z_i=0.2X_i.
\]

Therefore,

\[
\boxed{Z_i\sim Pa(4,0.2)}.
\]

Hence,

\[
\boxed{
S_I\sim\mathcal{CP}(10,Pa(4,0.8))
}
\]

and

\[
\boxed{
S_R\sim\mathcal{CP}(10,Pa(4,0.2))
}.
\]
\end{frame}

\begin{frame}{Solution: Moments of the Original Severity}
\protect\phantomsection\label{solution-moments-of-the-original-severity}
For

\[
X\sim Pa(4,1),
\]

the mean is

\[
E[X]=\frac{\lambda}{\alpha-1}
=\frac{1}{4-1}
=\frac{1}{3}.
\]

The second moment is

\[
E[X^2]
=
\frac{\Gamma(\alpha-2)\Gamma(3)}
{\Gamma(\alpha)}
\lambda^2.
\]

Thus, \(E[X^2]=\frac{\Gamma(2)\Gamma(3)}{\Gamma(4)}=\frac{1}{3}.\)
\end{frame}

\begin{frame}{Solution: Original Aggregate Moments}
\protect\phantomsection\label{solution-original-aggregate-moments}
For a compound Poisson model,

\[
E[S]=\lambda E[X]
\]

and

\[
\operatorname{Var}(S)=\lambda E[X^2].
\]

Here \(\lambda=10\), so

\[
E[S]
=
10\left(\frac{1}{3}\right)
=
\boxed{\frac{10}{3}}.
\]
\end{frame}

\begin{frame}{Solution: Original Aggregate Variance}
\protect\phantomsection\label{solution-original-aggregate-variance}
Using

\[
\operatorname{Var}(S)=\lambda E[X^2],
\]

we obtain

\[
\operatorname{Var}(S)
=
10\left(\frac{1}{3}\right).
\]

Therefore,

\[
\boxed{\operatorname{Var}(S)=\frac{10}{3}}.
\]
\end{frame}

\begin{frame}{Solution: Insurer's Mean}
\protect\phantomsection\label{solution-insurers-mean}
Since

\[
S_I=0.8S,
\]

we can use the scaling property of expectation:

\[
E[S_I]
=
0.8E[S].
\]

Therefore,

\[
E[S_I]
=
0.8\left(\frac{10}{3}\right)
=
\boxed{\frac{8}{3}}.
\]
\end{frame}

\begin{frame}{Solution: Insurer's Variance}
\protect\phantomsection\label{solution-insurers-variance}
For a constant \(c\),

\[
\operatorname{Var}(cS)=c^2\operatorname{Var}(S).
\]

Therefore,

\[
\operatorname{Var}(S_I)
=
(0.8)^2\operatorname{Var}(S).
\]

Hence,

\[
\operatorname{Var}(S_I)
=
(0.8)^2\left(\frac{10}{3}\right)
=
\boxed{\frac{32}{15}}.
\]
\end{frame}

\begin{frame}{Solution: Reinsurer's Mean}
\protect\phantomsection\label{solution-reinsurers-mean}
Similarly,

\[
S_R=0.2S.
\]

Therefore,

\[
E[S_R]
=
0.2E[S]
=
0.2\left(\frac{10}{3}\right).
\]

Thus,

\[
\boxed{E[S_R]=\frac{2}{3}}.
\]
\end{frame}

\begin{frame}{Solution: Reinsurer's Variance}
\protect\phantomsection\label{solution-reinsurers-variance}
Using the variance scaling property,

\[
\operatorname{Var}(S_R)
=
(0.2)^2\operatorname{Var}(S).
\]

Therefore,

\[
\operatorname{Var}(S_R)
=
(0.2)^2\left(\frac{10}{3}\right)
=
\boxed{\frac{2}{15}}.
\]
\end{frame}

\begin{frame}{Solution: Compare the Variances}
\protect\phantomsection\label{solution-compare-the-variances}
The sum of the two individual variances is

\[
\operatorname{Var}(S_I)+\operatorname{Var}(S_R)
=
\frac{32}{15}+\frac{2}{15}
=
\boxed{\frac{34}{15}}.
\]

The original aggregate variance is

\[
\operatorname{Var}(S)=\frac{10}{3}
=\frac{50}{15}.
\]

Therefore,

\[
\boxed{
\frac{34}{15}<\frac{50}{15}
}.
\]
\end{frame}

\begin{frame}{Solution: Why Are They Different?}
\protect\phantomsection\label{solution-why-are-they-different}
We know that

\[
S=S_I+S_R.
\]

Thus,

\[
\operatorname{Var}(S)
=
\operatorname{Var}(S_I)
+
\operatorname{Var}(S_R)
+
2\operatorname{Cov}(S_I,S_R).
\]

The covariance term is important because

\[
S_I=0.8S,
\qquad
S_R=0.2S.
\]
\end{frame}

\begin{frame}{Solution: The Covariance}
\protect\phantomsection\label{solution-the-covariance}
We have

\[
\operatorname{Cov}(S_I,S_R)
=
\operatorname{Cov}(0.8S,0.2S).
\]

Therefore,

\[
\operatorname{Cov}(S_I,S_R)
=
0.8(0.2)\operatorname{Var}(S).
\]

Hence,

\[
\operatorname{Cov}(S_I,S_R)
=
0.16\left(\frac{10}{3}\right)
=
\boxed{\frac{8}{15}}.
\]

So the variance decomposition becomes
\(\frac{32}{15}+\frac{2}{15}+2\left(\frac{8}{15}\right)=\frac{50}{15}=\frac{10}{3}.\)
\end{frame}

\begin{frame}{Key Takeaways}
\protect\phantomsection\label{key-takeaways}
Under proportional reinsurance,

\[
\boxed{S_I=\alpha S,\qquad S_R=(1-\alpha)S}.
\]

Therefore,

\[
E[S_I]=\alpha E[S],
\qquad
E[S_R]=(1-\alpha)E[S].
\]

and

\[
\operatorname{Var}(S_I)=\alpha^2\operatorname{Var}(S),
\qquad
\operatorname{Var}(S_R)=(1-\alpha)^2\operatorname{Var}(S).
\]

The insurer and reinsurer payments are \textbf{dependent}, so

\[
\operatorname{Var}(S_I)+\operatorname{Var}(S_R)
\neq
\operatorname{Var}(S).
\]

\begin{block}{Compound Binomial Distribution}
\protect\phantomsection\label{compound-binomial-distribution}
\begin{block}{From Individual Policies to Aggregate Loss}
\protect\phantomsection\label{from-individual-policies-to-aggregate-loss-1}
Insurance companies are often interested not only in the claim amount
arising from a single policy, but in the \textbf{total amount of claims
generated by an entire portfolio} during a specified period.

Consider a portfolio of \(n\) independent and identical policies.
Suppose that, during the period under consideration, each policy can
produce \textbf{at most one claim}.

For each policy, there are therefore two possible outcomes:

\begin{itemize}
\tightlist
\item
  the policy produces no claim;
\item
  the policy produces one claim.
\end{itemize}

Let \(p\) denote the probability that a particular policy produces a
claim. Then the probability of no claim is \(1-p\).

If \(N\) denotes the total number of claims in the portfolio, then

\[
N\sim\operatorname{Binomial}(n,p).
\]

The binomial distribution is therefore a natural frequency model when
there is a \textbf{fixed number of potential claim events} and each
potential event can occur at most once.

The aggregate claim amount, however, depends not only on how many claims
occur but also on the amount of each claim. This leads to the compound
binomial model.
\end{block}

\begin{block}{Why the Binomial Distribution?}
\protect\phantomsection\label{why-the-binomial-distribution-1}
The binomial frequency model is appropriate when the following
assumptions are reasonable:

\begin{enumerate}
\tightlist
\item
  There are a fixed number \(n\) of policies or potential claim events.
\item
  Each policy can produce at most one claim during the period.
\item
  Each policy has the same probability \(p\) of producing a claim.
\item
  Claim events for different policies are independent.
\end{enumerate}

Under these assumptions,

\[
N\sim\operatorname{Binomial}(n,p),
\]

with probability mass function

\[
P(N=k)
=
\binom{n}{k}p^k(1-p)^{n-k},
\qquad k=0,1,\ldots,n.
\]

An important feature of this model is that the number of claims has a
fixed upper bound:

\[
0\leq N\leq n.
\]

For example, if there are \(1{,}000\) policies, there cannot be more
than \(1{,}000\) claims under the assumption that each policy can
produce at most one claim.
\end{block}

\begin{block}{Frequency and Severity}
\protect\phantomsection\label{frequency-and-severity}
The random variable \(N\) describes the \textbf{frequency} of claims:
how many claims occur.

To describe the total amount paid, we also need a \textbf{severity}
model.

Let

\[
X_i
\]

denote the amount of the \(i\)th claim. Assume that

\[
X_1,X_2,\ldots
\]

are i.i.d. and independent of \(N\).

The aggregate claim amount is then

\[
S=\sum_{i=1}^{N}X_i.
\]

The number of terms in this sum is random because \(N\) is random.

For example:

\begin{itemize}
\tightlist
\item
  if \(N=0\), then \(S=0\);
\item
  if \(N=1\), then \(S=X_1\);
\item
  if \(N=2\), then \(S=X_1+X_2\);
\item
  if \(N=3\), then \(S=X_1+X_2+X_3\).
\end{itemize}

Thus, the aggregate loss is a \textbf{random sum}.

This gives the basic structure of the collective risk model:

\[
\boxed{
\text{Frequency}
\quad+\quad
\text{Severity}
\quad\longrightarrow\quad
\text{Aggregate loss}.
}
\]

The frequency component determines \textbf{how many claims occur}, while
the severity component determines \textbf{how much each claim costs}.
\end{block}
\end{block}

\begin{block}{The Compound Binomial Model}
\protect\phantomsection\label{the-compound-binomial-model-1}
\begin{block}{Definition}
\protect\phantomsection\label{definition}
Suppose that

\[
N\sim\operatorname{Binomial}(n,p),
\]

and that the claim amounts

\[
X_1,X_2,\ldots
\]

are i.i.d. and independent of \(N\).

The aggregate claim amount

\[
\boxed{
S=\sum_{i=1}^{N}X_i
}
\]

is said to have a \textbf{compound binomial distribution}.

We may write

\[
\boxed{
S\sim\mathcal{CB}(n,p,F_X),
}
\]

where \(F_X\) denotes the distribution of an individual claim amount.

The three quantities \((n,p,F_X)\) therefore determine the model:

\begin{itemize}
\tightlist
\item
  \(n\) describes the number of potential claim events;
\item
  \(p\) describes the probability of a claim for each policy;
\item
  \(F_X\) describes the distribution of the amount of an individual
  claim.
\end{itemize}
\end{block}

\begin{block}{Conditional Interpretation}
\protect\phantomsection\label{conditional-interpretation}
A useful way to understand the compound binomial distribution is to
condition on the number of claims.

If

\[
N=k,
\]

then the aggregate claim amount is

\[
S=X_1+\cdots+X_k.
\]

Therefore, conditional on \(N=k\), the aggregate loss is the sum of
\(k\) independent severity variables.

In particular,

\[
S\mid N=k
\overset{d}{=}
X_1+\cdots+X_k.
\]

Consequently, the distribution of \(S\) can be viewed as a mixture of
the distributions corresponding to different possible values of \(N\).

This is the origin of the word \textbf{compound}: the aggregate
distribution is constructed by combining the frequency distribution with
the severity distribution.
\end{block}

\begin{block}{Probability of Zero Claims}
\protect\phantomsection\label{probability-of-zero-claims}
The binomial frequency model also gives an immediate expression for the
probability of no claims:

\[
P(N=0)
=
(1-p)^n.
\]

If claim amounts are strictly positive, then

\[
S=0
\quad\Longleftrightarrow\quad
N=0.
\]

Therefore,

\[
\boxed{
P(S=0)=(1-p)^n.
}
\]

This is an important difference between an aggregate loss distribution
and an individual claim-size distribution. The aggregate loss has a
\textbf{point mass at zero} because the portfolio may experience no
claims during the period.
\end{block}

\begin{block}{Moments of the Aggregate Loss}
\protect\phantomsection\label{moments-of-the-aggregate-loss}
\end{block}

\begin{block}{Mean}
\protect\phantomsection\label{mean}
For a random sum

\[
S=\sum_{i=1}^{N}X_i,
\]

where \(N\) is independent of the severity variables, the expected value
is

\[
E[S]=E[N]E[X].
\]

For the binomial frequency distribution,

\[
E[N]=np.
\]

Let

\[
m_1=E[X].
\]

Then

\[
\begin{aligned}
E[S]
&=E[N]E[X]\\
&=(np)m_1.
\end{aligned}
\]

Hence,

\[
\boxed{
E[S]=npm_1.
}
\]

This result has a particularly simple interpretation:

\[
\boxed{
E[S]
=
\underbrace{np}_{\text{expected number of claims}}
\underbrace{m_1}_{\text{expected claim amount}}.
}
\]

Thus, the expected aggregate loss is the expected number of claims
multiplied by the expected amount of an individual claim.

For example, suppose

\[
n=10{,}000,
\qquad
p=0.02,
\]

and

\[
E[X]=50{,}000.
\]

Then the expected number of claims is

\[
E[N]=10{,}000(0.02)=200,
\]

and consequently

\[
E[S]
=
200(50{,}000)
=
10{,}000{,}000.
\]

The expected aggregate loss is therefore \(10\) million.

The calculation illustrates an important modelling principle:
\textbf{the mean aggregate loss can be understood through the two
separate components of the model}.
\end{block}

\begin{block}{Variance}
\protect\phantomsection\label{variance}
The variance of a random sum is

\[
\operatorname{Var}(S)
=
E[N]\operatorname{Var}(X)
+
\operatorname{Var}(N)\{E[X]\}^2.
\]

For the binomial distribution,

\[
E[N]=np
\]

and

\[
\operatorname{Var}(N)=np(1-p).
\]

Let

\[
m_1=E[X]
\]

and

\[
m_2=E[X^2].
\]

Since

\[
\operatorname{Var}(X)=m_2-m_1^2,
\]

we obtain

\[
\begin{aligned}
\operatorname{Var}(S)
&=
np(m_2-m_1^2)
+
np(1-p)m_1^2.
\end{aligned}
\]

Expanding,

\[
\begin{aligned}
\operatorname{Var}(S)
&=
npm_2-npm_1^2
+
np(1-p)m_1^2\\
&=
npm_2-npm_1^2
+npm_1^2-np^2m_1^2\\
&=
npm_2-np^2m_1^2.
\end{aligned}
\]

Therefore,

\[
\boxed{
\operatorname{Var}(S)
=
npm_2-np^2m_1^2.
}
\]

Equivalently,

\[
\boxed{
\operatorname{Var}(S)
=
np\operatorname{Var}(X)
+
np(1-p)\{E[X]\}^2.
}
\]

The latter expression provides an important interpretation.

The first component,

\[
np\operatorname{Var}(X),
\]

represents uncertainty arising from \textbf{variation in claim sizes}.

The second component,

\[
np(1-p)\{E[X]\}^2,
\]

represents uncertainty arising from \textbf{variation in the number of
claims}.

Thus,

\[
\boxed{
\text{Aggregate uncertainty}
=
\text{severity uncertainty}
+
\text{frequency uncertainty}.
}
\]

This decomposition is one of the most useful ways to interpret the
variance of a collective risk model.
\end{block}

\begin{block}{Third Central Moment and Coefficient of Skewness}
\protect\phantomsection\label{third-central-moment-and-coefficient-of-skewness}
The mean and variance describe the centre and spread of the aggregate
loss distribution. However, insurance loss distributions are often
asymmetric, with a relatively long right tail.

The \textbf{coefficient of skewness} measures this asymmetry.

For a random variable \(Y\), define

\[
\gamma_1(Y)
=
\frac{E[(Y-E[Y])^3]}
{\{\operatorname{Var}(Y)\}^{3/2}}.
\]

For the severity distribution, define the raw moments

\[
m_1=E[X],
\qquad
m_2=E[X^2],
\qquad
m_3=E[X^3].
\]

For the compound binomial model, the third central moment of the
aggregate loss is

\[
\boxed{
\mu_3(S)
=
np\left[
m_3-3pm_1m_2+2p^2m_1^3
\right].
}
\]

Therefore, the coefficient of skewness is

\[
\boxed{
\gamma_1(S)
=
\frac{
np\left[
m_3-3pm_1m_2+2p^2m_1^3
\right]
}{
\left(
npm_2-np^2m_1^2
\right)^{3/2}
}.
}
\]

A positive value of \(\gamma_1(S)\) indicates right-skewness.

This is common for insurance aggregate losses because unusually large
claims can produce aggregate losses far above the mean.

The skewness therefore gives information that cannot be obtained from
the mean and variance alone.
\end{block}

\begin{block}{A Useful Interpretation of Skewness}
\protect\phantomsection\label{a-useful-interpretation-of-skewness}
The third central moment can also be written as

\[
\begin{aligned}
\mu_3(S)
&=
np\mu_{3,X}
+
3np(1-p)\mu_X\sigma_X^2\\
&\quad
+
np(1-p)(1-2p)\mu_X^3,
\end{aligned}
\]

where

\[
\mu_X=E[X],
\]

\[
\sigma_X^2=\operatorname{Var}(X),
\]

and

\[
\mu_{3,X}
=
E[(X-\mu_X)^3].
\]

This expression shows that aggregate skewness is influenced by both:

\begin{itemize}
\tightlist
\item
  the skewness of the individual claim amounts;
\item
  the randomness in the number of claims.
\end{itemize}

Therefore, even if individual claims were not particularly skewed,
randomness in claim frequency could still contribute to aggregate
skewness.
\end{block}

\begin{block}{The MGF of the Compound Binomial Distribution}
\protect\phantomsection\label{the-mgf-of-the-compound-binomial-distribution}
The moment generating function provides another way to construct and
analyse the compound binomial distribution.

Let

\[
M_X(t)=E[e^{tX}]
\]

be the MGF of the individual claim amount.

Conditional on \(N=k\),

\[
S=X_1+\cdots+X_k.
\]

Since the \(X_i\) are independent,

\[
\begin{aligned}
M_{S\mid N=k}(t)
&=
E[e^{t(X_1+\cdots+X_k)}]\\
&=
\prod_{i=1}^{k}E[e^{tX_i}]\\
&=
[M_X(t)]^k.
\end{aligned}
\]

We now average over the possible values of \(N\):

\[
\begin{aligned}
M_S(t)
&=
E[M_X(t)^N].
\end{aligned}
\]

The probability generating function of a binomial random variable is

\[
G_N(z)
=
(1-p+pz)^n.
\]

Setting

\[
z=M_X(t)
\]

gives

\[
\boxed{
M_S(t)
=
[1-p+pM_X(t)]^n.
}
\]

Writing

\[
q=1-p,
\]

we may also write

\[
\boxed{
M_S(t)
=
[q+pM_X(t)]^n.
}
\]

This formula is the compound binomial counterpart of the general
random-sum MGF relationship.
\end{block}
\end{block}

\begin{block}{Worked Example: Compound Binomial with Gamma Severity}
\protect\phantomsection\label{worked-example-compound-binomial-with-gamma-severity-1}
Consider a portfolio with

\[
n=100
\]

policies, where each policy has probability

\[
p=0.01
\]

of producing a claim.

Thus,

\[
N\sim\operatorname{Binomial}(100,0.01).
\]

Suppose that, independently of \(N\),

\[
X_i\overset{\text{i.i.d.}}{\sim}\operatorname{Gamma}(20,0.1),
\]

where the Gamma distribution uses the \textbf{shape-rate
parameterisation}.

Let

\[
S=\sum_{i=1}^{N}X_i.
\]

We will determine:

\begin{enumerate}
\tightlist
\item
  the MGF of \(S\);
\item
  \(E[S]\);
\item
  \(\operatorname{Var}(S)\);
\item
  the coefficient of skewness;
\item
  \(P(S=0)\).
\end{enumerate}

\begin{block}{Moments of the Gamma Severity}
\protect\phantomsection\label{moments-of-the-gamma-severity}
For

\[
X\sim\operatorname{Gamma}(\alpha,\lambda),
\]

under the shape-rate parameterisation,

\[
E[X]=\frac{\alpha}{\lambda}
\]

and

\[
\operatorname{Var}(X)=\frac{\alpha}{\lambda^2}.
\]

Therefore,

\[
E[X]
=
\frac{20}{0.1}
=
200.
\]

Thus,

\[
m_1=200.
\]

The variance is

\[
\operatorname{Var}(X)
=
\frac{20}{0.1^2}
=
2{,}000.
\]

Hence,

\[
\begin{aligned}
m_2
&=E[X^2]\\
&=\operatorname{Var}(X)+\{E[X]\}^2\\
&=2{,}000+200^2\\
&=42{,}000.
\end{aligned}
\]

For the third raw moment of a Gamma random variable,

\[
E[X^3]
=
\frac{\alpha(\alpha+1)(\alpha+2)}
{\lambda^3}.
\]

Thus,

\[
m_3
=
\frac{20(21)(22)}{0.1^3}
=
9{,}240{,}000.
\]

Therefore,

\[
\boxed{
m_1=200,\qquad
m_2=42{,}000,\qquad
m_3=9{,}240{,}000.
}
\]
\end{block}

\begin{block}{MGF of the Aggregate Loss}
\protect\phantomsection\label{mgf-of-the-aggregate-loss-1}
For the Gamma severity,

\[
M_X(t)
=
\left(
\frac{0.1}{0.1-t}
\right)^{20},
\qquad t<0.1.
\]

Using the compound binomial MGF,

\[
M_S(t)
=
[1-p+pM_X(t)]^n.
\]

Substituting \(n=100\) and \(p=0.01\) gives

\[
\boxed{
M_S(t)
=
\left[
0.99+
0.01
\left(
\frac{0.1}{0.1-t}
\right)^{20}
\right]^{100}.
}
\]

This single expression describes the distribution of the aggregate claim
amount.
\end{block}

\begin{block}{Expected Aggregate Loss}
\protect\phantomsection\label{expected-aggregate-loss-3}
The expected number of claims is

\[
E[N]=np=100(0.01)=1.
\]

Therefore,

\[
\begin{aligned}
E[S]
&=npm_1\\
&=100(0.01)(200)\\
&=200.
\end{aligned}
\]

Hence,

\[
\boxed{
E[S]=200.
}
\]

The interpretation is straightforward: on average, one claim occurs and
each claim has expected amount \(200\).
\end{block}

\begin{block}{Variance of the Aggregate Loss}
\protect\phantomsection\label{variance-of-the-aggregate-loss}
Using

\[
\operatorname{Var}(S)
=
npm_2-np^2m_1^2,
\]

we obtain

\[
\begin{aligned}
\operatorname{Var}(S)
&=
100(0.01)(42{,}000)
-
100(0.01)^2(200)^2\\
&=
42{,}000-400\\
&=
41{,}600.
\end{aligned}
\]

Thus,

\[
\boxed{
\operatorname{Var}(S)=41{,}600.
}
\]

The standard deviation is

\[
\operatorname{SD}(S)
=
\sqrt{41{,}600}
\approx203.96.
\]

The standard deviation is quite large relative to the mean because the
Gamma severity distribution allows substantial variation in claim
amounts.
\end{block}

\begin{block}{Coefficient of Skewness}
\protect\phantomsection\label{coefficient-of-skewness}
Using

\[
\gamma_1(S)
=
\frac{
np\left[
m_3-3pm_1m_2+2p^2m_1^3
\right]
}{
\left(
npm_2-np^2m_1^2
\right)^{3/2}
},
\]

we obtain

\[
\begin{aligned}
\gamma_1(S)
&=
\frac{
100(0.01)
\left[
9{,}240{,}000
-3(0.01)(200)(42{,}000)
+2(0.01)^2(200)^3
\right]
}{
(41{,}600)^{3/2}
}\\
&\approx1.060.
\end{aligned}
\]

Hence,

\[
\boxed{
\gamma_1(S)\approx1.06.
}
\]

The positive skewness indicates that the aggregate-loss distribution has
a longer right tail than left tail.
\end{block}

\begin{block}{Probability of Zero Aggregate Loss}
\protect\phantomsection\label{probability-of-zero-aggregate-loss-1}
Because the claim amounts are positive,

\[
S=0
\quad\Longleftrightarrow\quad
N=0.
\]

Therefore,

\[
\begin{aligned}
P(S=0)
&=P(N=0)\\
&=(1-p)^n\\
&=(0.99)^{100}\\
&\approx0.3660.
\end{aligned}
\]

Thus,

\[
\boxed{
P(S=0)\approx0.3660.
}
\]

There is approximately a \(36.6%
\) probability that the portfolio generates no claims during the period.
\end{block}
\end{block}

\begin{block}{From the Compound Binomial Model to the General Collective
Risk Model}
\protect\phantomsection\label{from-the-compound-binomial-model-to-the-general-collective-risk-model}
The compound binomial distribution is an important special case of the
general collective risk model

\[
S=\sum_{i=1}^{N}X_i.
\]

The general structure remains the same:

\[
\boxed{
S
=
\underbrace{\sum_{i=1}^{N}}_{\text{random number of claims}}
\underbrace{X_i}_{\text{claim severity}}.
}
\]

What changes from one compound model to another is the distribution
chosen for \(N\).

For the compound binomial model,

\[
N\sim\operatorname{Binomial}(n,p),
\]

so there is a fixed maximum number of claims:

\[
N\leq n.
\]

This makes the compound binomial model particularly suitable for
portfolios in which each individual policy can generate at most one
claim.

In contrast, if a policy can generate multiple claims during the period,
a different frequency model, such as the Poisson or negative binomial
distribution, may be more appropriate.

The compound binomial model therefore provides an important bridge
between \textbf{individual-policy modelling} and the broader
\textbf{collective risk model}.

\begin{block}{Summary}
\protect\phantomsection\label{summary-1}
The compound binomial model combines a binomial frequency distribution
with a severity distribution.

The frequency model is

\[
N\sim\operatorname{Binomial}(n,p),
\]

and the aggregate loss is

\[
S=\sum_{i=1}^{N}X_i.
\]

If

\[
m_k=E[X^k],
\]

then

\[
\boxed{
E[S]=npm_1
}
\]

and

\[
\boxed{
\operatorname{Var}(S)
=
npm_2-np^2m_1^2.
}
\]

The third central moment is

\[
\boxed{
\mu_3(S)
=
np\left[
m_3-3pm_1m_2+2p^2m_1^3
\right],
}
\]

giving the coefficient of skewness

\[
\boxed{
\gamma_1(S)
=
\frac{
np\left[
m_3-3pm_1m_2+2p^2m_1^3
\right]
}{
\left(
npm_2-np^2m_1^2
\right)^{3/2}
}.
}
\]

The MGF is

\[
\boxed{
M_S(t)
=
[1-p+pM_X(t)]^n.
}
\]

The central modelling idea is that the aggregate loss contains two
distinct sources of randomness:

\[
\boxed{
\text{Frequency uncertainty}
+
\text{Severity uncertainty}
=
\text{Aggregate loss uncertainty}.
}
\]

This same framework will be used for other compound distributions, with
the principal difference being the choice of the frequency distribution
\(N\).
\end{block}
\end{block}
\end{frame}




\end{document}
